\section{Data Preparation and Training-View Construction}
\label{app:dataset}

% 介绍对原始数据集的处理：
% 1.清洗数据，剔除了原始数据里面的不好的东西（具体是哪些）
% 2.修改数据，对原始数据里面的 prompt 做了哪些修改（具体改了哪些）
% 3.介绍最终的数据情况（哪些数据集，数量多少，怎么划分的，示例

\subsection{Preparing Data for RS and SO}

% 对作者提供的数据进行任务筛选和样本质量清洗。
% 原始数据包含 8 个域内任务维度，最终保留 7 个，
% 排除 \texttt{verifiability\_extraction}。
% 对四个 RevUtil 任务分别构造 4,800 条候选样本，
% 删除 human-hard 样本、测试集重叠样本和异常字符记录。
% 两个数据集的训练池均经过教师蒸馏，仅保留格式有效且教师分数
% 与金标签一致的轨迹（教师蒸馏的详细配置与所用 prompt 见
% 附录~\ref{app:teacher-distillation}）。测试集另外删除 coherence
% 的 2 条重复样本，以及 positioning_check 的 2 条训练集重叠样本。
% Related Work 和 RevUtil 两个数据集的最终规模及评分范围
% 分别汇总于表~\ref{tab:rw-data-size} 和表~\ref{tab:rev-util-data-size}。

The source contains eight in-domain task dimensions; task and sample-quality
filtering retained seven, excluding \texttt{verifiability\_extraction}. For
each of the four RevUtil tasks, a 4,800-example candidate pool was constructed
after removing human-hard examples, test-set overlaps, and malformed-character
records. Both datasets use training pools constructed through teacher distillation,
retaining only trajectories with valid formatting and teacher scores matching
the gold labels. Detailed distillation settings and prompts are provided in
Appendix~\ref{app:teacher-distillation}. 
Test-set cleanup additionally removed two duplicate
\texttt{coherence} examples and two training-overlap
\texttt{positioning\_check} examples.The final sizes and scoring ranges of the 
RevUtil and Related Work datasets are
summarized in Tables~\ref{tab:rw-data-size} and~\ref{tab:rev-util-data-size},
respectively.

% 从每个 prompt 中删除完整的 \texttt{[EXAMPLES]} 区块，
% 以及 system prompt 中关于提供示例的句子；原始 query、criteria、
% 评分规则和待评价文本保持不变。为最终记录加入稳定 ID 和监督元数据。
% 每个样本均构造成两种严格对齐的监督形式：rationale-and-score（RS）和 score-only（SO）。
% 两种形式共享样本、输入、金标签和数据划分，

Every prompt was stripped of the complete \texttt{[EXAMPLES]} block and the
corresponding system-prompt sentence, while the original query, criteria,
scoring rubric, and evaluated instance were preserved. Stable identifiers and
supervision metadata were added to the resulting records. Each example was 
instantiated in two strictly aligned supervision formats:
rationale-and-score (RS) and score-only (SO). The two formats share the same
examples, inputs, gold labels, and data splits; only their system instructions
and completions differ.

% 最终训练池使用 seed=20260720 和 0.1 的验证集比例进行划分。
% 表中的 Train / Val. 表示 optimizer training split 和 validation split。

\subsection{Constructing ReaSon Training Data}

% 两方法的第二阶段训练数据均由各自第一阶段模型（按各自训练目标、seed=42
% 训练得到）在相同推理配置下生成，生成分数统一替换为金标签。
% 七个任务覆盖三个 Related Work 数据集和四个 RevUtil 数据集，
% 共 13,680 条训练池样本，按 seed=20260720 划分为
% 12,312 条训练样本和 1,368 条验证样本。

For both methods, the second-stage training data are regenerated by the
first-stage model of each respective method (trained with its own training
objective, seed 42) under identical inference configurations, with generated
scores replaced by gold scores. The seven tasks cover three Related Work and
four RevUtil datasets, totaling 13,680 training-pool examples, split with
seed 20260720 into 12,312 training and 1,368 validation instances.

\subsection{Constructing PaIR Training Views}

% PaIR 覆盖 7 个任务和 13,680 条源样本，每条样本配对构造 DS 与 RF 两个视图，
% 两视图在同一 micro-batch 中出现，通过各自的 loss mask 区分。
% 训练时分别计算两视图 token 级交叉熵（prompt token 的 labels 设为 -100，
% 仅 completion token 参与监督），每个视图的 loss 在各自监督 token 上取平均，
% 总 loss 为 0.5·L_DS + 0.5·L_RF。
% 配对训练沿用 seed=20260720 和 0.1 validation ratio，得到 12,312 条
% optimizer-training 样本和 1,368 条 validation 样本；测试集不参与训练。

PaIR covers seven tasks and 13,680 source examples, each paired with a DS
view and an RF view that appear in the same micro-batch. During training,
token-level cross-entropy is computed over the two views separately via
per-view loss masks (prompt tokens are masked to \(-100\), so only completion
tokens contribute to supervision); each view's loss is averaged over its own
supervised tokens, and the total loss is
\(0.5\,\mathcal{L}_{\mathrm{DS}} + 0.5\,\mathcal{L}_{\mathrm{RF}}\).
Using seed 20260720 and a validation ratio of 0.1, the paired data yield
12,312 optimizer-training and 1,368 validation examples; test sets remain
outside training for matched evaluation.

% RevUtil 任务的最终数据规模。
\begin{table}[H]
\centering
\small
\begin{tabular}{lrrr}
\toprule
\textbf{Task} & Train / Val. & Test & Scoring \\
\midrule
Actionability & 1,609 / 179 & 1,000 & 1--5 \\
Grounding Spec. & 2,387 / 265 & 1,000 & 1--5 \\
Helpfulness & 2,051 / 228 & 1,000 & 1--5 \\
Verifiability & 1,643 / 182 & 788 & 1--5 \\
\midrule
\textbf{Total} & 7,690 / 854 & 3,788 & \textbf{12332} \\
\bottomrule
\end{tabular}
\caption{Final statistics and scoring ranges for the RevUtil tasks. 
The total scoring range sums the scores across the four tasks.}
\label{tab:rev-util-data-size}
\end{table}

% Related Work 任务的最终数据规模。
\begin{table}[H]
\centering
\small
\begin{tabular}{lrrr}
\toprule
\textbf{Task} & Train / Val. & Test & Scoring \\
\midrule
Coherence & 1,373 / 153 & 1,046 & 0--1 \\
Positioning Cons. & 2,399 / 267 & 603 & 0--1 \\
Positioning Type & 850 / 94 & 204 & 0--1 \\
\midrule
\textbf{Total} & 4,622 / 514 & 1,853 & \textbf{6,989} \\
\bottomrule
\end{tabular}
\caption{Final statistics and scoring ranges for the Related Work tasks. 
The total scoring range sums the scores across the three tasks.}
\label{tab:rw-data-size}
\end{table}
